% !TeX root = ../../Book2.tex
% 纲要标题：### 12.2 障碍余圈及其不变量
% 纲要位置：几何学计划纲要.md，第 1478 行。
% 正文状态：初稿；尚未正式审核。

% 纲要内容开始
% 
% * 障碍上链的构造
% * 余圈条件的证明
% * 改变初始延拓时的余边界
% * 障碍类消失与延拓存在
% * 差异类与上同调可表性
% 
% 纲要内容结束

\section{障碍余圈及其不变量}
\label{g2:sec:1202}

一个边界球面类判断一个胞腔能否填入。对整个空间，胞腔的边界之间还有相互关系，而重新选择已经填好的部分也会改变下一步的障碍。上链记录逐胞腔的值，余圈条件记录这些值之间的关系，除去余边界则记录允许改变上一层后的不变量。本节逐一证明这三件事。

\subsection{障碍上链的构造}
\label{g2:subsec:1202:01}

固定道路连通 CW 对 \((X,A)\)、Hurewicz 纤维化 \(p:E\to X\)，以及 \(A\) 上的指定截面 \(s_A\)。假定纤维非空且简单，记上一节构造的局部系统为 \(\Pi_n(p)\)。对某个 \(n\ge1\)，设已经选定延拓 \(s:X_n=A\cup X^n\to E\)。

\ProseParagraphBreak

给每个相对 \((n+1)\) 胞腔选定向及特征映射 \(\chi_e:D^{n+1}\to X_{n+1}\)。沿圆盘的径向收缩把边界截面送到中心纤维，便得到一个球面映射 \(a_e:S^n\to F_{\chi_e(0)}\)。源球面取圆盘的正向边界；它的同伦类就是这个胞腔上的延拓障碍。

\begin{definition}[障碍上链]\label{g2:def:obstruction-cochain}
	在上述条件下，将每个有向相对 \((n+1)\) 胞腔 \(e\) 赋值为 \([a_e]\in\pi_n(F_{\chi_e(0)})\)，得到\term[zhangai-shanglian]{障碍上链}[obstruction cochain]
	\[
		c_{n+1}(s)\in C^{n+1}(X,A;\Pi_n(p)),\qquad c_{n+1}(s)(e)=[a_e].
	\]
	在 \(A\) 内的胞腔上取零。若通过道路把系数搬到选定的基点纤维，则同时按第5章的局部系数约定记录运输。
\end{definition}

这一赋值包含三项选择，但它们的作用可分别追踪。反转胞腔定向使球面类变号；改变圆盘上的径向提升，只产生纤维内同伦；改变连接基点的道路，则按 \(\Pi_n(p)\) 的运输变换系数。纤维简单性保证同一纤维内的基点道路不再产生歧义。因此得到的是一个确定的局部系数上链，而不是把不同纤维中的元素直接相加。

\ProseParagraphBreak

由\cref{g2:prop:section-cell-extension}，保持 \(s\) 在 \(X_n\) 上的全部值时，延拓到 \(X_{n+1}\) 的充要条件是
\begin{equation}\label{g2:eq:obstruction-cochain-vanishing}
	c_{n+1}(s)=0.
\end{equation}
这是上链逐项为零的条件。它与稍后出现的上同调类为零有不同含义。

\begin{example}[带一个关系的二维复形]\label{g2:exm:obstruction-relation-cochain}
	令 \(X=S^1\cup_m e^2\)，二胞腔沿圆周的 \(m\) 次绕行附着，\(m\ge2\)。给定一次骨架映射 \(f:S^1\to S^1\)，其次数为 \(k\)。把它看成乘积纤维化的截面，则
	\[
		c_2(f)(e^2)=mk\in\pi_1(S^1)=\mathbb Z.
	\]
	因为 \(m\ne0\)，保持这个 \(f\) 延拓的条件是 \(k=0\)。另一方面，胞腔上链的余边界 \(C^1(X;\mathbb Z)\to C^2(X;\mathbb Z)\) 是乘 \(m\)，所以 \([c_2(f)]=0\) 对所有 \(k\) 都成立。允许改变一胞腔上的映射后，确实总能把它换为常值再延拓。这个例子先把两种零值判据的差别显示出来。
\end{example}

\subsection{余圈条件的证明}
\label{g2:subsec:1202:02}

虽然 \(s\) 尚未定义在 \((n+1)\) 胞腔内部，障碍值仍须满足 \((n+2)\) 胞腔给出的关系。关键是：整个特征圆盘可缩，它的拉回纤维化有截面；这个截面未必与已有边界数据相同，但足以证明障碍值在胞腔边界上的总和为零。

\begin{lemma}[障碍余圈]\label{g2:lem:obstruction-cocycle}
	设 \((X,A)\) 是道路连通 CW 对，\(p:E\to X\) 为纤维非空且简单的 Hurewicz 纤维化。对 \(n\ge1\) 及截面 \(s:X_n\to E\)，有
	\[
		\delta c_{n+1}(s)=0\quad\text{于 }C^{n+2}(X,A;\Pi_n(p)).
	\]
\end{lemma}
\begin{proof}
	先把逐胞腔的定义扩大到相对圆盘。给定带基点映射 \(\Phi:(D^{n+1},S^n)\to(X_{n+1},X_n)\)，用 \(s\Phi\) 作为边界提升，沿 \(\Phi\) 给出的径向同伦提升到一个纤维，记所得类为 \(\omega_s[\Phi]\)。相对同伦的提升说明它只依赖相对同伦类；将两个圆盘沿加法方向拼接，再提升同一拼接，说明 \(\omega_s\) 是群同态。若 \(\Phi\) 本身已有一个与指定边界相同的提升，则所得球面可填，\(\omega_s[\Phi]=0\)。

	为了同时处理运输，拉回到 \(X\) 的万有覆盖 \(\widetilde X\)。局部系统在那里成为常值系统；对每个骨架记其全部逆像为 \(\widetilde X_r\)。当 \(r\ge2\) 时，\(X_r\hookrightarrow X\) 在基本群上为同构，所以 \(\widetilde X_r\) 单连通。对 \(n\ge2\)，空间对 \((\widetilde X_{n+1},\widetilde X_n)\) 只附着 \((n+1)\) 胞腔，故为 \(n\) 连通；\cref{g2:thm:hurewicz-relative}给出
	\[
		\pi_{n+1}(\widetilde X_{n+1},\widetilde X_n)\xrightarrow{\cong}
		H_{n+1}(\widetilde X_{n+1},\widetilde X_n).
	\]
	特征圆盘在这张同构下成为相应的有向胞腔。因此 \(\omega_s\) 在相对同调上的表示，恰好是障碍上链的求值，而不是另选的一张同态。

	现在取一个 \((n+2)\) 胞腔的特征映射 \(g:D^{n+2}\to\widetilde X_{n+2}\)，并在边界上选基点。其胞腔边界由以下路径得到：先取 \(g|_{S^{n+1}}\) 的球面类，再把它送入 \(\pi_{n+1}(\widetilde X_{n+1},\widetilde X_n)\)，最后用上述 Hurewicz 同构。这与同调边界一致，因为 Hurewicz 映射与边界及包含相容，见\cref{g2:prop:hurewicz-naturality}。

	把纤维化拉回 \(D^{n+2}\)，可选一个在边界基点处等于 \(s\) 的截面。它给出 \(g|_{S^{n+1}}\) 的整体提升。将这个球面写成 \(D^{n+1}/S^n\) 后，所得相对圆盘在其边界恒取选定基点，提升在该边界也恒等于 \(s\) 的值。因此这一相对类被 \(\omega_s\) 送到零。换言之，\(c_{n+1}(s)\) 在任一 \((n+2)\) 胞腔的胞腔边界上取零。

	\(n=1\) 时，\(\widetilde X_1\) 未必单连通，需要把上述相对同伦到胞腔链的步骤直接说明。对二维相对圆盘，将它经过的二维胞腔表示为有向小圆盘，用细弧连接到边界基点；剩余部分映到 \(\widetilde X_1\)。沿细弧切开后，每个小圆盘贡献其胞腔障碍的正值或负值，剩余部分已有截面，贡献为零。重新粘合时，边上的道路成对反向出现；改变连接弧只把相应纤维环路共轭。由于 \(\pi_1(F)\) 交换，共轭消失，所得总值只依赖各二维胞腔的有向次数。这就证明 \(\omega_s\) 经由 \(H_2(\widetilde X_2,\widetilde X_1)\) 分解。上述小圆盘表示来自相对胞腔逼近及圆盘映射的次数分解，与第3章计算胞腔关联数所用的分解相同。

	对三胞腔的附着球面使用这一分解，其每条内部边出现两次、方向相反；也可沿前一段的圆盘提升，直接看出相应相对类的 \(\omega_s\) 为零。因此 \(n=1\) 同样得到 \(\delta c_2(s)=0\)。最后把覆盖中的计算降回 \(X\)。改变提升胞腔会同时改变其中心纤维的识别，恰好插入 \(\Pi_n(p)\) 的运输；第5章的局部系数胞腔余边界正是这一等变求值。因此结论对原来的局部系数成立。\(A\) 内的胞腔原已有截面，所有相应值为零，故上述论证也保留相对条件。
\end{proof}

由引理可定义
\begin{equation}\label{g2:eq:obstruction-class}
	o_{n+1}(s)=[c_{n+1}(s)]\in H^{n+1}(X,A;\Pi_n(p)).
\end{equation}
这称为当前截面的\term[zhangai-lei]{障碍类}[obstruction class]。上同调群描述障碍可能出现的位置；\(o_{n+1}(s)\) 是当前数据给出的具体元素。一个非零群完全可能包含等于零的障碍类。

\subsection{改变初始延拓时的余边界}
\label{g2:subsec:1202:03}

设 \(s,t\) 都定义在 \(X_n\)，并在 \(X_{n-1}\) 上相同。对一个 \(n\) 胞腔，将 \(t\) 的圆盘取正向、\(s\) 的圆盘取负向，沿共同边界粘成 \(S^n\)，再利用圆盘上的纤维同伦模型送入中心纤维。所得类记为
\[
	D_n(s,t)(e^n)\in\Pi_n(p)_{\chi_e(0)}.
\]
这些值组成\term[chayi-shanglian]{差异上链}[difference cochain] \(D_n(s,t)\in C^n(X,A;\Pi_n(p))\)。它记录把 \(s\) 改成 \(t\) 时，在每个 \(n\) 胞腔中实际加入的球面类。

\ProseParagraphBreak

若 \(s,t\) 只在 \(X_{n-1}\) 上由指定截面同伦 \(h\) 相连，则用 \(h\) 填侧面，将上下两张圆盘联成柱体的整个边界，同样得到 \(D_n(s,t;h)\)。为固定符号，柱体取先胞腔、后区间的乘积定向。若 \(H\) 是两端用 \(s,t\)、侧面用 \(h\) 的部分截面，则
\begin{equation}\label{g2:eq:difference-prism-sign}
	D_n(s,t;h)(e^n)=(-1)^n c_{n+1}(H)(e^n\times I).
\end{equation}
因为 \(\partial(e^n\times I)=\partial e^n\times I+(-1)^n(e^n\times\{1\}-e^n\times\{0\})\)，这一约定恰好使上端为 \(t\)、下端为 \(-s\)。

\begin{proposition}[差异与余边界]\label{g2:prop:obstruction-difference}
	设 \((X,A)\) 为道路连通 CW 对，\(p:E\to X\) 为纤维非空且简单的 Hurewicz 纤维化，\(n\ge1\)。若 \(s,t:X_n\to E\) 是在 \(A\) 上相同的截面，\(h\) 是它们在 \(X_{n-1}\) 上固定 \(A\) 的截面同伦，则
	\begin{equation}\label{g2:eq:obstruction-difference}
		\delta D_n(s,t;h)=c_{n+1}(t)-c_{n+1}(s).
	\end{equation}
	此外，给定 \(s\) 和任意 \(b\in C^n(X,A;\Pi_n(p))\)，都能选到在 \(X_{n-1}\) 上与 \(s\) 相同的 \(t\)，使 \(D_n(s,t)=b\)。
\end{proposition}
\begin{proof}
	令 \(W=X_{n+1}\)，在 \(W\times I\) 上拉回纤维化。部分截面 \(H\) 已在 \(X_n\times\partial I\)、\(X_{n-1}\times I\) 和 \(A\times I\) 上定义，这包含相对柱体的 \(n\) 骨架。由上一引理，其 \((n+1)\) 次障碍在每个 \((n+2)\) 胞腔 \(e^{n+1}\times I\) 的边界上取零。因此
	\[
		0=\sum_{e^n}[e^{n+1}:e^n]c_{n+1}(H)(e^n\times I)
		+(-1)^{n+1}\bigl(c_{n+1}(t)-c_{n+1}(s)\bigr)(e^{n+1}).
	\]
	有局部系数时，求和中同时沿各附着道路运输；两端的系数用区间方向的恒等运输识别。代入式\eqref{g2:eq:difference-prism-sign}并乘 \((-1)^n\)，便得到所需公式。

	为实现指定的 \(b\)，在每个相对 \(n\) 胞腔内部取一个小圆盘。在纤维模型中，把该圆盘的边界压成一点，插入代表 \(b(e^n)\) 的球面映射；原来通向小圆盘中心的值通过一条纤维道路与球面基点连接。等价地，这是保持外边界的挤压映射 \(D^n\to D^n\vee S^n\)，在第一因子沿用原映射，在第二因子加入所选球面类。简单性使连接道路不影响所得类，方向取为 \(t-s\)，于是差异正是 \(b(e^n)\)。对 \(n=1\)，此操作就是在原道路中插入指定环路。

	纤维同伦模型转换回原丛后，若外边界只相等到纤维同伦，按\cref{g2:prop:section-cell-extension}证明中的相对提升把它修正为实际相等。各胞腔内部的修改相互独立，在 \(X_{n-1}\) 及 \(A\) 上全不改变，故由 CW 弱拓扑粘合成 \(t\)。它满足 \(D_n(s,t)=b\)。
\end{proof}

因此，保持 \(X_{n-1}\) 上的数据时，\(c_{n+1}(s)\) 恰好可以改变一个任意余边界。障碍类正是这组修改下保持不变的量。普通映射的对应叙述见\cite[第18章第5节，Obstruction theory]{May1999}；本节的局部运输与柱体公式给出了截面版本的同一机制。

\subsection{障碍类消失与延拓存在}
\label{g2:subsec:1202:04}

现在可以把“修改后填入”的几何要求写成一个上同调判据。

\begin{theorem}[逐层延拓判据]\label{g2:thm:obstruction-extension}
	设 \((X,A)\) 为道路连通 CW 对，\(p:E\to X\) 为纤维非空且简单的 Hurewicz 纤维化，\(n\ge1\)，且 \(s:X_n\to E\) 为截面。则以下两条件等价：
	\begin{enumerate}
		\item \(o_{n+1}(s)=0\in H^{n+1}(X,A;\Pi_n(p))\)；
		\item \(s|_{X_{n-1}}\) 可保持不变地延拓为 \(X_{n+1}\) 上的截面。
	\end{enumerate}
	若要求保持原来的整个 \(s|_{X_n}\)，充要条件是上链 \(c_{n+1}(s)=0\)。
\end{theorem}
\begin{proof}
	若 \(o_{n+1}(s)=0\)，选 \(b\in C^n(X,A;\Pi_n(p))\) 使 \(\delta b=-c_{n+1}(s)\)。差异命题给出与 \(s\) 在 \(X_{n-1}\) 上相同的 \(t\)，满足
	\[
		c_{n+1}(t)=c_{n+1}(s)+\delta b=0.
	\]
	单胞腔判据把 \(t\) 延拓到 \(X_{n+1}\)。反过来，若有这样的延拓，记它在 \(X_n\) 上的限制为 \(t\)，则 \(c_{n+1}(t)=0\)，而差异公式使 \(c_{n+1}(s)=-\delta D_n(s,t)\)，故障碍类为零。最后一句直接是\cref{g2:eq:obstruction-cochain-vanishing}。
\end{proof}

\begin{theorem}[首障碍]\label{g2:thm:primary-obstruction}
	设 \((X,A)\) 为道路连通 CW 对，\(p:E\to X\) 为纤维非空、简单且 \((n-1)\) 连通的 Hurewicz 纤维化，\(n\ge1\)，并给定截面 \(s_A\)。则 \(s_A\) 总可延到 \(X_n\)，且
	\[
		o_{n+1}(p,s_A)\in H^{n+1}(X,A;\Pi_n(p))
	\]
	与所选 \(X_n\) 上的延拓无关。它为零当且仅当 \(s_A\) 可延到 \(X_{n+1}\)。这个\term[shouzhangai]{首障碍}[primary obstruction]对保持给定截面的 CW 对拉回自然。若相对胞腔维数均不超过 \(n+1\)，它为零便等价于全空间上的截面存在。
\end{theorem}
\begin{proof}
	非空性允许选择新顶点的值，道路连通性允许填入新一胞腔；\(\pi_j(F)=0\)（\(0<j<n\)）再使维数至 \(n\) 的胞腔都可以填入。故存在所需初始延拓。

	取两张这样的截面 \(s,t\)。在 \(X_{n-1}\) 上构造固定 \(A\) 的截面同伦，只需处理 \(k\) 胞腔乘区间，\(k\le n-1\)；其边界问题取值于 \(\pi_k(F)\)。当 \(k=0\) 时由道路连通性填入，其余均由低阶同伦群消失填入。于是得到前一命题所需的 \(h\)，差异公式证明两张障碍上圈同属一个上同调类。

	若这个类为零，逐层延拓判据给出延到 \(X_{n+1}\) 的截面；反之，取实际延拓的限制来计算，障碍上链本身为零。对胞腔映射的拉回，在每个相对胞腔上，边界球面复合及运输与原构造相容；将复合映射分解成有向胞腔贡献，就得到局部系数上链上的拉回公式。一般映射先相对胞腔逼近，拉回纤维化在基底同伦下纤维同伦等价，故得到相同的障碍类。这证明自然性。最后，在所述维数范围内 \(X_{n+1}=X\)，所以无需再处理更高胞腔。
\end{proof}

纤维较低同伦群的消失在这里承担两个不同任务：它使初始截面存在，也使更低骨架上的初始选择同伦唯一。缺少第二项时，第 \(n+1\) 次障碍类仍对固定的低层数据有意义，却可能随着更低层的选择改变。

\begin{corollary}[由障碍群消失得到延拓]\label{g2:cor:obstruction-groups-vanish}
	设 \((X,A)\) 为道路连通 CW 对，\(p:E\to X\) 为纤维非空且简单的 Hurewicz 纤维化。若
	\[
		H^{q+1}(X,A;\Pi_q(p))=0\qquad(q\ge1),
	\]
	则每个 \(A\) 上的截面都可延拓到 \(X\)。对有限维相对 CW 对，只需核对实际出现的次数。
\end{corollary}
\begin{proof}
	先延拓到 \(X_1\)。在第 \(q\) 步，障碍群为零使当前障碍类为零，所以可以保持 \(X_{q-1}\) 不变，修改 \(q\) 胞腔上的值，再延拓到 \(X_{q+1}\)。一个固定的 \(k\) 胞腔只会在第 \(k\) 步最后修改一次，以后所有步骤都保留它。取这些最终值，便在每个特征圆盘上得到连续且边界相容的截面；CW 弱拓扑使它们定义整个 \(X\) 上的连续截面。
\end{proof}

\subsection{差异类与上同调可表性}
\label{g2:subsec:1202:05}

截面存在以后，还可以问它是否同伦唯一。这一次要填的是两张截面之间的柱体，因此障碍所在的上同调次数降低一。

\begin{proposition}[同伦的差异类]\label{g2:prop:section-difference-class}
	设 \((X,A)\)、\(p:E\to X\) 满足上一推论的空间与纤维条件，\(s,t\) 为在 \(A\) 上相同的全局截面。给定它们在 \(X_{n-1}\) 上固定 \(A\) 的截面同伦 \(h\)，\(n\ge1\)，则 \(D_n(s,t;h)\) 为余圈。其类
	\[
		\Delta_n(s,t;h)\in H^n(X,A;\Pi_n(p))
	\]
	为零，当且仅当 \(h|_{X_{n-2}}\) 可延拓为 \(s|_{X_n}\) 与 \(t|_{X_n}\) 之间的截面同伦。\(n=1\) 时 \(X_{-1}=A\)，所以只固定 \(A\) 上的同伦。
\end{proposition}
\begin{proof}
	两张全局截面的障碍上链均为零，所以\cref{g2:eq:obstruction-difference}给出 \(\delta D_n=0\)。在空间对
	\[
		\bigl(X\times I,\ X\times\partial I\cup A\times I\bigr)
	\]
	上考虑给定的两端与部分同伦。相对胞腔为 \(e^k\times(0,1)\)，维数为 \(k+1\)。相对链的边界去掉两端后只有 \(\partial e^k\times I\)，所以其局部系数上同调自然识别为 \(H^{j-1}(X,A;\Pi_n(p))\)。这就是柱体相对方向使次数降低一的原因。

	用式\eqref{g2:eq:difference-prism-sign}识别同伦障碍，逐层延拓判据便说明：其类为零，恰好允许修改内部的 \(n\) 维胞腔，即 \(e^{n-1}\times I\) 上的同伦，同时固定更低柱体与全部两端，再填入 \(e^n\times I\)。这正是结论所述的延拓，而不是要求保留原来的整个 \(h|_{X_{n-1}}\)。
\end{proof}

如果纤维 \((n-1)\) 连通，上述 \(h\) 在所需低次中总能找到。其不同选择在更低的柱体上又可同伦，因为相关同伦群仍在消失范围内；差异公式于是使 \(\Delta_n\) 的上同调类与这些低次选择无关。接下来两种情形中，这个差异类已经记录全部自由度。

\begin{theorem}[截面同伦类的参数]\label{g2:thm:section-torsor}
	设 \((X,A)\) 为道路连通 CW 对，\(p:E\to X\) 为具有非空、简单、\((n-1)\) 连通纤维的 Hurewicz 纤维化，\(n\ge1\)，并假设指定截面 \(s_A\) 至少有一个全局延拓。以下任一条件成立时，延拓的相对截面同伦类形成 \(H^n(X,A;\Pi_n(p))\) 上的主齐性集：
	\begin{enumerate}
		\item 相对胞腔维数均不超过 \(n\)；
		\item 纤维的全部正阶同伦群除 \(\pi_n\) 外均为零。
	\end{enumerate}
	选定一个延拓 \(s_0\) 后，\([t]\mapsto\Delta_n(s_0,t)\) 给出与该上同调群的双射。
\end{theorem}
\begin{proof}
	低阶同伦群消失使各截面在 \(X_{n-1}\) 上可由相对同伦相连。选定这样的同伦后，逐胞腔球面的拼接给出
	\[
		\Delta_n(s,t)+\Delta_n(t,u)=\Delta_n(s,u),\qquad
		\Delta_n(s,s)=0.
	\]
	改变辅助低次同伦只加一个余边界：将两次同伦再放到一个参数区间上比较，在维数低于 \(n\) 的地方由纤维连通度填充，剩余柱体边界正是差异公式。因而这些等式确实是上同调中的等式。

	若 \(\Delta_n(s,t)=0\)，上一命题给出 \(X_n\) 上的相对同伦。条件 I 下已经完成全空间；条件 II 下，继续填入更高的 \(e^k\times I\) 所遇系数 \(\pi_k(F)\) 均为零，所以同伦延拓到全部 \(X\)。反过来，全局同伦当然使每一级差异障碍为零。这证明所列映射单射。

	给定任意上同调类，取余圈 \(b\in C^n(X,A;\Pi_n(p))\) 表示它。在 \(s_0\) 的各 \(n\) 胞腔中插入 \(b\) 的球面代表，得到 \(t_n\)，并有 \(D_n(s_0,t_n)=b\)。条件 I 下它已经是全局截面；条件 II 下，由 \(\delta b=0\) 及 \(c_{n+1}(s_0)=0\)，得到 \(c_{n+1}(t_n)=0\)，先延到 \(X_{n+1}\)，再利用更高同伦群消失继续延拓。所得截面给出指定的差异类，证明满射。

	因此上同调群通过加入差异类自由且传递地作用于延拓同伦类。这样的集合称为\term[zhuqixingji]{主齐性集}[torsor]：任意两个元素之间有唯一的群元素作为差，但在选择 \(s_0\) 前没有一个自然的零元。
\end{proof}

\begin{theorem}[上同调的可表性]\label{g2:thm:em-representability}
	设 \(G\) 是阿贝尔群，\(n\ge1\)，\((X,A)\) 为 CW 对。取带基点 CW 模型 \(K(G,n)\)，并令 \(\iota_n\in H^n(K(G,n);G)\) 为在 \(\pi_n\cong G\) 上对应恒等映射的类。则有自然双射
	\[
		[(X,A),(K(G,n),*)]\xrightarrow{\cong}H^n(X,A;G),\qquad [f]\longmapsto f^*\iota_n.
	\]
	左端是把 \(A\) 映到基点的映射在固定 \(A\) 的同伦下的类；无相对条件时得到 \([X,K(G,n)]\cong H^n(X;G)\)。
\end{theorem}
\begin{proof}
	第9章已经构造了 \(K(G,n)\)。\(n\ge2\) 时，Hurewicz 定理给出 \(H_n(K;\mathbb Z)\cong G\) 以及低正次同调消失，故万有系数定理把 \(H^n(K;G)\) 识别为 \(\operatorname{Hom}(G,G)\)。\(n=1\) 时，用一次 Hurewicz 定理及 \(G\) 交换得到相同识别；\(H_0=\mathbb Z\) 不产生 Ext 项。因此 \(\iota_n\) 定义明确。

	对乘积纤维化 \(X\times K(G,n)\to X\)，常值映射是指定边界值的一个延拓。纤维为 \((n-1)\) 连通简单空间，且只剩一个正阶同伦群，所以上一定理的条件 II 给出由 \(H^n(X,A;G)\) 参数化的全部映射同伦类。\(X\) 不连通时逐分支应用即可。

	还要验证所得参数确实是 \(f^*\iota_n\)。先把 \(f\) 在 \(X_{n-1}\) 上同伦为常值，并由同伦延拓性质将这段同伦延到全空间。对每个相对 \(n\) 胞腔，此时 \(f\) 压去边界后给出 \(S^n\to K(G,n)\)。差异上链在该胞腔上的值就是此球面代表的 \(G\) 中元素；而 \(\iota_n\) 对这个球面基本类的求值也正是该元素，因为它对应 \(G\) 上的恒等映射。因此两个胞腔上圈相同，给出所列拉回公式。公式与空间对映射的复合相容，所以双射自然。
\end{proof}

这一结论把一个抽象上同调类变成一张实际映射。下一节正要利用它：把一个阶段的粘合类表示为到 \(K(G,n+1)\) 的映射，再取该映射的同伦纤维，便能把相应的 \(G\) 放回所需次数。关于可表性、Postnikov 阶段与障碍问题之间的关系，可参看\cite[第4.3节，Connections with Cohomology]{Hatcher2002}。
% 纲要条目排为编号小节标题，并进入目录与书签；说明段仍保留于纲要注释中。
